\subsection*{MT-ART11. The tame integral class is a central-parity invariant}
Let $N/K$ be a finite Galois extension of number fields, tamely
ramified at every finite place, with group $G$. All module classes
in this statement are over $\mathbb Z[G]$, after restriction of
scalars from $\mathcal O_K$. Put
\[
 U_{N/K}=[\mathcal O_N]-[\mathcal O_K[G]]
       \in\operatorname{Cl}(\mathbb Z[G]),\qquad
 m_\chi=\operatorname{ord}_{s=1/2}L(s,\chi,N/K)\in\mathbb Z.
                                                        \tag{ART11.1}
\]
Orders here are meromorphic orders, so a negative value denotes
a pole. Let $R_G^s$ be the virtual symplectic-character group and
$T(\chi)=\chi+\overline\chi$. The classical root-number map
\[
 t:\operatorname{Hom}_{\Omega_{\mathbb Q}}
       (R_G^s/T(R_G),\{\pm1\})
          \longrightarrow\operatorname{Cl}(\mathbb Z[G])
                                                        \tag{ART11.2}
\]
is additive after writing its sign-group law multiplicatively.
Then
\[
 U_{N/K}=t\bigl(\chi\longmapsto(-1)^{m_\chi}\bigr),
 \qquad 2U_{N/K}=0.                                    \tag{ART11.3}
\]
In particular,
\[
 \zeta_N(1/2)\ne0\ \Longrightarrow\ U_{N/K}=0,
 \qquad
 U_{N/K}\ne0\ \Longrightarrow\
        \operatorname{ord}_{s=1/2}\zeta_N(s)\ge2.          \tag{ART11.4}
\]
Neither implication uses Artin holomorphy or GRH. The first
means stable freeness as a $\mathbb Z[G]$-module of rank
$[K:\mathbb Q]$; it does not by itself assert a relative normal
integral basis over $\mathcal O_K[G]$.

\paragraph{Classical input.}
Taylor, \emph{On Fr\"ohlich's conjecture for rings of integers
of tame extensions}, Invent.\ Math.\ 63 (1981), 41--79,
Theorem~1, p.~41, gives $U=t(W)$ and $2U=0$.
Cassou-Nogu\`es--Taylor, Ann.\ Inst.\ Fourier 33(2) (1983),
1--17, (2.2)--(2.5), pp.~6--7, and Corollary~A, p.~10,
fix (ART11.2) and the relative normalization in (ART11.1).
The original Taylor theorem page was visually inspected;
the 1983 paper's full main proof, p.~11, was read with its
explicit upstream Gauss-sum theorems retained as citations.

\begin{proof}
The completed Artin function has a meromorphic functional
equation. If $\chi$ is symplectic it is self-dual, so
\[
 \Lambda(1/2+z,\chi)=W(\chi)\Lambda(1/2-z,\chi).
\]
Its archimedean gamma and conductor factors are holomorphic
and nonzero at $1/2$. Write its nonzero Laurent expansion as
$a z^{m_\chi}+O(z^{m_\chi+1})$, where $a\ne0$ and
$m_\chi$ may be negative. Comparing the two leading
coefficients gives $1=W(\chi)(-1)^{m_\chi}$.
The same calculation extends multiplicatively to virtual
symplectic characters. Thus the parity character in
(ART11.3) is exactly the classical root-number character;
in particular it has the Galois invariance and the
$T(R_G)$-quotient property needed in (ART11.2).
Taylor's theorem gives (ART11.3).

Here are full details of the nonvanishing step, so that no
unproved holomorphy assertion enters (ART11.4). Set
\[
             \Theta=\sum_{\chi\in\operatorname{Irr}(G)}
                         m_\chi\chi.
\]
Suppose $\zeta_N(1/2)\ne0$. For a cyclic subgroup
$C\le G$, put $F=N^C$. The abelian factorization is
\[
       \zeta_N(s)=\prod_{\eta\in\widehat C}
                       L(s,\eta,N/F).
\]
Every factor is holomorphic at $1/2$, including the trivial
factor $\zeta_F$. Their nonnegative orders sum to zero,
so each order is zero. Artin induction and Frobenius
reciprocity identify these orders with
$\langle\operatorname{Res}^G_C\Theta,\eta\rangle_C$.
Hence $\operatorname{Res}^G_C\Theta=0$ for every cyclic
$C$. Every $g\in G$ lies in $\langle g\rangle$, so
$\Theta(g)=0$ for every $g$, and irreducible-character
independence gives $m_\chi=0$ for every $\chi$.
The root-number character is therefore trivial and
$U_{N/K}=0$. This also proves that nonvanishing of
$\zeta_N$ excludes central Artin poles in this argument.

For the contrapositive, $\zeta_N$ is holomorphic at $1/2$.
Its completed functional equation has sign $+1$ and
nonzero archimedean factors there. A central zero
therefore has positive even order, proving the lower
bound two in (ART11.4).
\end{proof}

If $G$ has no irreducible symplectic characters, its
symplectic quotient $R_G^s/T(R_G)$ is zero and $U_{N/K}=0$.
For odd $|G|$, Taylor's Corollary, p.~41, with complete
proof p.~42, strengthens this to actual freeness over
$\mathbb Z[G]$, using the Eichler condition and Jacobinski
cancellation. This cancellation step is part of that
classical corollary. There is no general freeness claim
from a vanishing stable class in the preceding proof.

For use beside the positivity gate MT-ART02, there is the
following precise refinement. If all the central Artin
orders are nonnegative, then
\[
 \operatorname{ord}_{1/2}\zeta_N
   =\sum_{\chi\in\operatorname{Irr}(G)}\chi(1)m_\chi
   \ \ge\!
   \sum_{\substack{\chi\ {
m irreducible\ symplectic}\\W(\chi)=-1}}
                         \chi(1).                       \tag{ART11.5}
\]
Indeed each selected order is a nonnegative odd integer
and therefore at least one; every other term is
nonnegative. Thus an integral sign obstruction gives a
degree-weighted lower bound once the central Heilbronn
matrix is positive. Without that extra positivity, the
unconditional conclusion remains (ART11.4).

The placement is the locally free integral $K_0$ of the
finite Galois cover, in contrast with its regular rational
representation after the normal-basis theorem. It extends
the concrete integral distinction of MT-ART10 and
AG-LEFSCH-LATTICE/AG-LEFSCH-DEFECT. In the inherited
cote~16, scans~14--18, no.~7 integral-ring discussion,
it supplies an exact sign-valued invariant that survives
in the lattice but is absent from the rational module.
The associated Artin representation lies in the
SGA~6, V, \S2 and \S7, especially (7.1.1),
representation-$\lambda$-ring crosswalk fixed in AG-LAMBDA.
These are programme placements, not an assertion that
SGA~6 or the archive contains Taylor's arithmetic theorem.

\subsection*{MT-ART12. The hermitian trace class recovers every central parity}
Keep the tame extension and the right-module convention.
Equip $\mathcal O_N$ with its $G$-valued trace form
and the comparison module with its standard trace form:
\[
 \begin{split}
 h_N(m,n)&=\sum_{\gamma\in G}
                \operatorname{Tr}_{N/\mathbb Q}(m n^\gamma)
                                \gamma^{-1},\\
 h_0(m,n)&=\operatorname{Tr}_{K/\mathbb Q}(\overline m n),
       \qquad m,n\in K[G],\\
 X_{N/K}&=[\mathcal O_N,h_N]-[\mathcal O_K[G],h_0]
                   \in K_0H(\mathbb Z[G]).
 \end{split}                                             \tag{ART12.1}
\]
The bar in $h_0$ is the group-ring involution
$\gamma\mapsto\gamma^{-1}$, and the field trace is
applied coefficientwise. Here $K_0H$ is the Grothendieck
group for these locally free hermitian modules and
orthogonal sum. Its discriminant has the canonical
sign extraction $\theta$ of Cassou-Nogu\`es--Taylor,
(2.9), p.~8. Their main theorem, p.~10, and
Corollary~B, p.~11, give
\[
 \theta_\chi(X_{N/K})=W(\chi),\qquad
 \eta_\chi(X_{N/K})=
       \bigl(\operatorname{N}\mathfrak f(\chi)\bigr)^{1/2}
                                      W_\infty(\chi)
                                                        \tag{ART12.2}
\]
for every symplectic $\chi$. These exact normalizations
are (3.6)--(3.9), p.~9; the original page was inspected,
including the square root of the conductor and the bar
in $h_0$.

The consequence for the present spectral programme is
\[
       (-1)^{\operatorname{ord}_{1/2}L(s,\chi)}
                         =\theta_\chi(X_{N/K}).           \tag{ART12.3}
\]
Thus the hermitian trace class determines each
symplectic central parity individually, including
meromorphic orders, without GRH or Artin holomorphy.
The underlying locally free class is obtained by the
forgetful map and gives $t(W)$ as in MT-ART11; inversion
of $t$ is not needed in (ART12.3).

\begin{proof}
Apply the complete Laurent-series parity calculation
in MT-ART11 to each symplectic character and substitute
the classical first equality in (ART12.2). This proves
(ART12.3) with no positivity premise. If two tame
extensions over the same base have identified groups
and the same class $X$ under that identification, the
homomorphisms $\theta_\chi$ take the same values, so
all their symplectic central orders have the same
parity. Applying the forgetful map gives equal
underlying stable classes, by (ART11.3). Conversely,
this proof makes no injectivity assumption on $t$;
stable-class data alone have not been substituted
for the hermitian class.
\end{proof}

The definition of the reduced hermitian group in
the cited source uses extension to
$\overline{\mathbb Q}[G]$; the source's p.~5 and p.~10
were visually checked to retain that algebraic closure.
No isometry over $\mathbb Q[G]$ is silently inferred.
The construction supplies a precise hermitian refinement
of the integral $K_0$ comparison in MT-ART11 and of
the trace programme at cote~1. It gives root signs
and central parities, while nonnegative multiplicities
remain the distinct Heilbronn-positivity gate. The
extra conditional lower bound (ART11.5) states exactly
what that positivity would add.
